Fix the following Hilbert calculus for theoremhood in this language.  A
derivation is a finite sequence of formulas in which each line is a uniform
substitution instance of a classically valid propositional formula, an
instance of
one of the three schemas
\[
\begin{aligned}
  \mathsf X(A\to B)&\to(\mathsf XA\to\mathsf XB),\\
  \mathsf X\neg A&\leftrightarrow\neg\mathsf XA,\\
  A\mathbin{\mathsf U}B
    &\leftrightarrow
      B\lor(A\land\mathsf X(A\mathbin{\mathsf U}B)),
\end{aligned}
\]
or follows from earlier lines by one of the rules
\[
  \frac{A\qquad A\to B}{B}\;\textsc{MP},
  \qquad
  \frac{A}{\mathsf XA}\;\textsc{XNec},
  \qquad
  \frac{B\lor(A\land\mathsf X C)\to C}
       {A\mathbin{\mathsf U}B\to C}\;\textsc{UInd}.
\]
All metavariables in these schemas and rules range over temporal formulas.
Write \(\derives_{\mathsf{LTL}}A\) when such a no-premise derivation ends in
\(A\).  This is a syntactic theoremhood relation, distinct from temporal
validity, and it is introduced because the next theorem compares those two
standings.  Since derivations have no open premises, both necessitation and the
induction rule apply only to theorems.

